% Chapter 4 - Assembly.

\guidechapter{PROGRAMMING IN ASSEMBLY}{%
  \item The register and shared-variable conventions
  \item Relocating the SDK's RAM and zero page
  \item Every service entry point by example
  \item Raw request blocks with optional spans
}

\section{The Calling Convention}

Include \cmd{ultimate.inc} and call the unprefixed entry points. Unless an
entry below says otherwise, a service returns an \cmd{ULTIMATE\_*} result in
\cmd{A} and may change \cmd{A}, \cmd{X}, \cmd{Y} and the flags. A byte argument
uses \cmd{A}; a pointer argument uses low byte in \cmd{A} and high byte in
\cmd{X}; multi-byte service arguments live in the exported \cmd{ult\_*}
variables.

The examples use these two ordinary ca65 macros to keep pointer and word setup
visible without repeating four instructions each time:

\begin{listingbox}
\hspace*{2em}.include "ultimate.inc"\
\
.macro setptr location, value\
\hspace*{2em}lda \#<value\
\hspace*{2em}sta location\
\hspace*{2em}lda \#>value\
\hspace*{2em}sta location + 1\
.endmacro\
\
.macro setword location, value\
\hspace*{2em}lda \#<value\
\hspace*{2em}sta location\
\hspace*{2em}lda \#>value\
\hspace*{2em}sta location + 1\
.endmacro
\end{listingbox}

\begin{outputbox}
setptr ult_buf, text stores the address of text in ult_buf.
\end{outputbox}

\cmd{setptr} and \cmd{setword} emit the same bytes; their different names make
the purpose clear at each call site. Result branches always compare \cmd{A}
immediately, before another instruction replaces it.

Assembly output boxes show the returned register or memory values from a
representative successful run. They omit screen-printing routines so the SDK
operation itself stays visible.

With the C64 character map, lowercase source letters assemble to the uppercase
byte values shared by PETSCII and ASCII. Paths, filenames and HTTP text
therefore appear in lowercase source below but reach the firmware in uppercase.

\evidence{\cmd{bindings/asm/ultimate.inc}; \cmd{docs/asm-abi.md}}

\section{Place The SDK's Working Memory}

The SDK's zero-page block starts at \cmd{UCI\_ZP}, which defaults to
\cmd{\$FB}. Its size is \cmd{UCI\_ZP\_SIZE}. The normal variable block occupies
\cmd{UCI\_VARS\_SIZE} bytes in the BSS segment, which is where the linker puts
variables that have no starting value. Define both bases when assembling every
SDK source to choose fixed locations instead:

\begin{listingbox}
ca65 -D UCI\_VARS=\$CF00 -D UCI\_ZP=\$A3 uci\_core.s
\end{listingbox}

\begin{outputbox}
UCI variables begin at \$CF00; zero-page workspace begins at \$A3.
\end{outputbox}

With \cmd{UCI\_VARS} defined, the core emits no BSS; its variables become
equates from that address. The SDK installs no interrupt handler and is not
re-entrant, so an interrupt routine must not call it while foreground code is
inside a service.

\evidence{\cmd{bindings/asm/uci\_protocol.inc};
\cmd{src/uci/uci\_zp.inc}; \cmd{docs/asm-abi.md}}

\Needspace{30\baselineskip}
\section{Initialize, Detect And Describe The Machine}

\begin{listingbox}
\hspace*{2em}jsr ultimate\_init\
\hspace*{2em}cmp \#ULTIMATE\_OK\
\hspace*{2em}bne failed\
\
\hspace*{2em}jsr ultimate\_available\
\hspace*{2em}beq unavailable\
\
\hspace*{2em}lda \#<capabilities\
\hspace*{2em}ldx \#>capabilities\
\hspace*{2em}jsr ultimate\_detect\
\hspace*{2em}cmp \#ULTIMATE\_OK\
\hspace*{2em}bne failed\
\
capabilities: .res 4
\end{listingbox}

\begin{outputbox}
A = ULTIMATE_OK\\
capabilities records the detected targets.
\end{outputbox}

\cmd{ultimate\_init} performs startup. \cmd{ultimate\_available} returns a
boolean rather than a result code. \cmd{ultimate\_detect} takes the address of
the four-byte capability block directly in \cmd{A}/\cmd{X}.

\Needspace{20\baselineskip}
\subsection*{Identify a target with and without the optional length output}

\begin{listingbox}
\hspace*{2em}setptr ult\_buf, text\
\hspace*{2em}setword ult\_buflen, 64\
\hspace*{2em}setptr ult\_outlen, text\_length\
\hspace*{2em}lda \#UCI\_TARGET\_DOS1\
\hspace*{2em}jsr ultimate\_identify\
\
\hspace*{2em}setword ult\_outlen, 0\
\hspace*{2em}lda \#UCI\_TARGET\_DOS1\
\hspace*{2em}jsr ultimate\_identify\
\
text:        .res 64\
text\_length: .res 2
\end{listingbox}

\begin{outputbox}
text = "ULTIMATE-II DOS V1.2"\\
text_length = 20 with the pointer; no length is stored with zero.
\end{outputbox}

Both calls terminate the text in \cmd{ult\_buf}. The first stores its length
through \cmd{ult\_outlen}; zero disables that optional output.

\Needspace{25\baselineskip}
\subsection*{Read the model with and without the optional length output}

\begin{listingbox}
\hspace*{2em}setptr ult\_buf, text\
\hspace*{2em}setword ult\_buflen, 64\
\hspace*{2em}setptr ult\_outlen, text\_length\
\hspace*{2em}jsr ultimate\_get\_model\
\
\hspace*{2em}setword ult\_outlen, 0\
\hspace*{2em}jsr ultimate\_get\_model
\end{listingbox}

\begin{outputbox}
text = "Ultimate 64 Elite"\\
text_length is optional in the same way as IDENTIFY.
\end{outputbox}

\cmd{ultimate\_get\_model} uses the same buffer contract but needs no target
argument. Its underlying \cmd{CTRL\_CMD\_GET\_HWINFO} command is marked
deprecated by the firmware documentation; this existing wrapper remains for
compatibility.

\Needspace{32\baselineskip}
\subsection*{Read SID addresses through deprecated HWINFO}

\begin{listingbox}
\hspace*{2em}lda \#<sid\_info\
\hspace*{2em}ldx \#>sid\_info\
\hspace*{2em}jsr ultimate\_legacy\_get\_sid\_info\
\hspace*{2em}cmp \#ULTIMATE\_OK\
\hspace*{2em}bne failed\
\
\hspace*{2em}lda sid\_info + ULTIMATE\_SID\_INFO\_COUNT\
\hspace*{2em}; iterate this many ULTIMATE_SID_RECORD_SIZE-byte records\
\hspace*{2em}ldx \#ULTIMATE\_SID\_INFO\_RECORDS\
\hspace*{2em}lda sid\_info + ULTIMATE\_SID\_PRIMARY\_ADDRESS,x\
\hspace*{2em}sta first\_sid\_address\
\hspace*{2em}lda sid\_info + ULTIMATE\_SID\_PRIMARY\_ADDRESS + 1,x\
\hspace*{2em}sta first\_sid\_address + 1\
\
sid\_info:          .res ULTIMATE\_SID\_INFO\_SIZE\
first\_sid\_address: .res 2
\end{listingbox}

\begin{outputbox}
sid_info count = 4\\
primary addresses = \$D400, \$D400, \$D400, \$D400
\end{outputbox}

Records begin at \cmd{ULTIMATE\_SID\_INFO\_RECORDS}. Within each record,
\cmd{ULTIMATE\_SID\_PRIMARY\_ADDRESS} and
\cmd{ULTIMATE\_SID\_SECONDARY\_ADDRESS} name two little-endian words, and
\cmd{ULTIMATE\_SID\_TYPE} names the raw type byte. The type is not a 6581 or
8580 identifier. Firmware still implements this query, but marks HWINFO
deprecated and publishes no replacement C64-side UCI command. The legacy name
is deliberate; handle \cmd{ULTIMATE\_ERR\_NOT\_SUPPORTED} and read no record
unless the call returned \cmd{ULTIMATE\_OK}.

\evidence{\cmd{bindings/asm/ultimate.inc}; \cmd{src/uci/ultimate.s};
\cmd{software/io/command\_interface/control\_target.cc}}

\Needspace{15\baselineskip}
\subsection*{Turn a result into text}

\begin{listingbox}
\hspace*{2em}lda error\_code\
\hspace*{2em}jsr ultimate\_strerror\
\hspace*{2em}sta message\_ptr\
\hspace*{2em}stx message\_ptr + 1\
\
error\_code: .res 1\
message\_ptr: .res 2
\end{listingbox}

\begin{outputbox}
message_ptr points to the PETSCII text for error_code.
\end{outputbox}

\cmd{ultimate\_strerror} returns the PETSCII message pointer in
\cmd{A}/\cmd{X}; it does not return a result code.

\evidence{\cmd{docs/asm-abi.md}; \cmd{src/uci/ultimate.s};
\cmd{src/uci/ultimate\_strerror.s}}

\Needspace{32\baselineskip}
\section{Read And Change The Palette}

\begin{listingbox}
\hspace*{2em}lda \#<palette\
\hspace*{2em}ldx \#>palette\
\hspace*{2em}jsr ultimate\_palette\_get\
\
\hspace*{2em}lda \#<palette\
\hspace*{2em}ldx \#>palette\
\hspace*{2em}jsr ultimate\_palette\_set\
\
\hspace*{2em}lda \#6\
\hspace*{2em}sta ult\_color\
\hspace*{2em}lda \#255\
\hspace*{2em}sta ult\_color + 1\
\hspace*{2em}lda \#0\
\hspace*{2em}sta ult\_color + 2\
\hspace*{2em}sta ult\_color + 3\
\hspace*{2em}jsr ultimate\_palette\_set\_color\
\
\hspace*{2em}jsr ultimate\_palette\_reset\
\
palette: .res UCI\_PALETTE\_BYTES
\end{listingbox}

\begin{outputbox}
GET fills 48 bytes; SET COLOUR 6 TO RED and RESET return ULTIMATE_OK.
\end{outputbox}

Get and set take a direct pointer in \cmd{A}/\cmd{X}. The single-colour entry
reads index, red, green and blue from the four consecutive \cmd{ult\_color}
bytes. Reset restores the built-in live palette.
Firmware without these commands returns
\cmd{ULTIMATE\_ERR\_NOT\_SUPPORTED}; branch on \cmd{A} before using the
48-byte result from \cmd{ultimate\_palette\_get}.

\evidence{\cmd{docs/asm-abi.md}; \cmd{src/uci/palette.s}}

\section{Use Turbo And Badline Control}

\begin{listingbox}
\hspace*{2em}jsr ultimate\_turbo\_available\
\hspace*{2em}beq no\_turbo\
\
\hspace*{2em}jsr ultimate\_turbo\_get\
\hspace*{2em}sta old\_speed\
\
\hspace*{2em}lda \#U64\_SPEED\_4MHZ\
\hspace*{2em}jsr ultimate\_turbo\_set\
\
\hspace*{2em}lda \#0\
\hspace*{2em}jsr ultimate\_turbo\_badlines\
\hspace*{2em}lda \#1\
\hspace*{2em}jsr ultimate\_turbo\_badlines\
\
\hspace*{2em}lda old\_speed\
\hspace*{2em}jsr ultimate\_turbo\_set\
\
old\_speed: .res 1
\end{listingbox}

\begin{outputbox}
old_speed = previous speed index\\
SET 4 MHZ, THEN RESTORE SPEED AND BADLINES: ULTIMATE_OK
\end{outputbox}

The available and get calls return values directly. Set returns a result. The
two badline calls show both forms explicitly: zero disables them; non-zero
restores them. The final set restores the speed read at entry.

\evidence{\cmd{docs/asm-abi.md}; \cmd{src/uci/turbo.s}}

\section{Change And Read The Current Directory}

\begin{listingbox}
\hspace*{2em}setptr ult\_buf, path\
\hspace*{2em}jsr ultimate\_chdir\
\
\hspace*{2em}setptr ult\_buf, path\_buffer\
\hspace*{2em}setword ult\_buflen, 128\
\hspace*{2em}setptr ult\_outlen, path\_length\
\hspace*{2em}jsr ultimate\_getpath\
\
\hspace*{2em}setword ult\_outlen, 0\
\hspace*{2em}jsr ultimate\_getpath\
\
path:        .byte "/usb1/games", 0\
path\_buffer: .res 128\
path\_length: .res 2
\end{listingbox}

\begin{outputbox}
path_buffer = "/USB1/GAMES"\\
path_length = 11 when the optional pointer is enabled.
\end{outputbox}

\cmd{ultimate\_chdir} consumes the terminated path through \cmd{ult\_buf}.
The two getpath calls show the optional length pointer enabled and disabled.

\subsection*{Read every directory entry}

\begin{listingbox}
\hspace*{2em}jsr ultimate\_opendir\
\hspace*{2em}cmp \#ULTIMATE\_OK\
\hspace*{2em}bne failed\
\
next\_entry:\
\hspace*{2em}setptr ult\_buf, name\
\hspace*{2em}setword ult\_buflen, 64\
\hspace*{2em}jsr ultimate\_readdir\
\hspace*{2em}cmp \#ULTIMATE\_END\
\hspace*{2em}beq directory\_done\
\hspace*{2em}cmp \#ULTIMATE\_OK\
\hspace*{2em}bne failed\
\hspace*{2em}lda ult\_attrib\
\hspace*{2em}; use name and its DOS\_ATTR\_* bits here\
\hspace*{2em}jmp next\_entry\
\
name: .res 64
\end{listingbox}

\begin{outputbox}
name receives one entry per call; A becomes ULTIMATE_END after the last.
\end{outputbox}

Each readdir releases one reply block and writes attributes to
\cmd{ult\_attrib}. Issue no other command inside the walk. Call
\cmd{uci\_abort} if leaving before \cmd{ULTIMATE\_END}.

\subsection*{Make a directory, and return to the configured one}

\begin{listingbox}
\hspace*{2em}setptr ult\_buf, dir\_name\
\hspace*{2em}jsr ultimate\_mkdir\
\
\hspace*{2em}jsr ultimate\_home\
\
dir\_name: .byte "saves", 0
\end{listingbox}

\begin{outputbox}
SAVES is created; A = ULTIMATE_OK from each call.
\end{outputbox}

\cmd{ultimate\_mkdir} takes the terminated name through \cmd{ult\_buf} and
measures it itself. \cmd{ultimate\_home} takes nothing at all.

\evidence{\cmd{docs/asm-abi.md}; \cmd{src/uci/dos.s}}

\section{Open, Seek, Read, Write And Close Files}

\begin{listingbox}
\hspace*{2em}setptr ult\_buf, read\_name\
\hspace*{2em}lda \#DOS\_FA\_READ\
\hspace*{2em}jsr ultimate\_open\
\hspace*{2em}cmp \#ULTIMATE\_OK\
\hspace*{2em}bne failed\
\
\hspace*{2em}lda \#0\
\hspace*{2em}sta ult\_num\
\hspace*{2em}sta ult\_num + 1\
\hspace*{2em}sta ult\_num + 2\
\hspace*{2em}sta ult\_num + 3\
\hspace*{2em}jsr ultimate\_seek\
\
\hspace*{2em}setptr ult\_buf, file\_buffer\
\hspace*{2em}setword ult\_buflen, 128\
\hspace*{2em}setptr ult\_outlen, bytes\_read\
\hspace*{2em}jsr ultimate\_read\
\hspace*{2em}jsr ultimate\_close\
\
read\_name:   .byte "readme.txt", 0\
file\_buffer: .res 128\
bytes\_read:  .res 2
\end{listingbox}

\begin{outputbox}
file_buffer receives up to 128 bytes; bytes_read receives the actual count.
\end{outputbox}

Open establishes the firmware's current file. Seek selects its absolute
position. Read stores bytes and, when \cmd{ult\_outlen} is non-zero, their
count. Close releases that same firmware-side file.

\Needspace{13\baselineskip}
The count output is optional:

\begin{listingbox}
\hspace*{2em}; With a readable file already open:\
\hspace*{2em}setptr ult\_buf, file\_buffer\
\hspace*{2em}setword ult\_buflen, 128\
\hspace*{2em}setword ult\_outlen, 0\
\hspace*{2em}jsr ultimate\_read
\end{listingbox}

\begin{outputbox}
file_buffer receives bytes; no byte count is stored.
\end{outputbox}

\subsection*{Create, write and delete a file}

\begin{listingbox}
\hspace*{2em}setptr ult\_buf, write\_name\
\hspace*{2em}lda \#DOS\_FA\_CREATE\_ALWAYS | DOS\_FA\_WRITE\
\hspace*{2em}jsr ultimate\_open\
\hspace*{2em}cmp \#ULTIMATE\_OK\
\hspace*{2em}bne failed\
\
\hspace*{2em}setptr ult\_buf, message\
\hspace*{2em}setword ult\_buflen, message\_end - message\
\hspace*{2em}jsr ultimate\_write\
\hspace*{2em}jsr ultimate\_close\
\
\hspace*{2em}setptr ult\_buf, write\_name\
\hspace*{2em}jsr ultimate\_delete\
\
write\_name: .byte "result.txt", 0\
message:    .byte "ok", 13\
message\_end:
\end{listingbox}

\begin{outputbox}
RESULT.TXT: 3 bytes written, closed and deleted; A = ULTIMATE_OK.
\end{outputbox}

The open mask creates or replaces the file. Write consumes
\cmd{ult\_buflen} bytes from \cmd{ult\_buf}; delete consumes the terminated
name after the file is closed.

\evidence{\cmd{docs/asm-abi.md}; \cmd{src/uci/dos.s}}

\Needspace{28\baselineskip}
\section{Ask What A File Is}

\begin{listingbox}
\hspace*{2em}setptr ult\_buf, stat\_name\
\hspace*{2em}setptr ult\_arg2, file\_info\
\hspace*{2em}jsr ultimate\_stat\
\hspace*{2em}cmp \#ULTIMATE\_OK\
\hspace*{2em}bne failed\
\
\hspace*{2em}lda file\_info + DOS\_INFO\_SIZE\
\hspace*{2em}sta file\_size\
\hspace*{2em}lda file\_info + DOS\_INFO\_SIZE + 1\
\hspace*{2em}sta file\_size + 1\
\hspace*{2em}lda file\_info + DOS\_INFO\_ATTRIB\
\hspace*{2em}and \#DOS\_ATTR\_DIR\
\hspace*{2em}bne is\_a\_directory\
\
stat\_name: .byte "readme.txt", 0\
file\_info: .res ULTIMATE\_FILEINFO\_SIZE\
file\_size: .res 2
\end{listingbox}

\begin{outputbox}
file_info holds the reply; the name inside it is NUL-terminated.
\end{outputbox}

\cmd{ult\_buf} is the name and \cmd{ult\_arg2} is the block the reply is
received into. That block must be \cmd{ULTIMATE\_FILEINFO\_SIZE} bytes.
\cmd{ult\_buflen} is not an input here: \cmd{ultimate\_stat} measures the name
itself and overwrites \cmd{ult\_buflen} doing it.

\cmd{ultimate\_fstat} asks the same question about the file
\cmd{ultimate\_open} left open, so it takes only the block:

\begin{listingbox}
\hspace*{2em}setptr ult\_arg2, file\_info\
\hspace*{2em}jsr ultimate\_fstat
\end{listingbox}

\begin{outputbox}
file_info describes the open file; ult_buf is not read.
\end{outputbox}

The \cmd{DOS\_INFO\_*} equates in \cmd{uci\_protocol.inc} name the fields
inside the block:

\begin{reftable}{lp{0.60\linewidth}}
\cmd{DOS\_INFO\_SIZE} & \cmd{+\$00}, size in bytes, 32-bit, low byte first \\
\cmd{DOS\_INFO\_DATE} & \cmd{+\$04}, year less 1980 in bits 15--9, month 8--5, day 4--0 \\
\cmd{DOS\_INFO\_TIME} & \cmd{+\$06}, hour 15--11, minute 10--5, two-second units 4--0 \\
\cmd{DOS\_INFO\_EXTENSION} & \cmd{+\$08}, three bytes, space padded \\
\cmd{DOS\_INFO\_ATTRIB} & \cmd{+\$0B}, the \cmd{DOS\_ATTR\_*} bits \\
\cmd{DOS\_INFO\_NAME} & \cmd{+\$0C}, the name, terminated by the SDK \\
\end{reftable}

Date and time are the packed forms used by FAT, the File Allocation Table
filesystem the Ultimate's storage uses, and they arrive unconverted.

\evidence{\cmd{bindings/asm/ultimate.inc}; \cmd{src/uci/dosinfo.s};
\cmd{bindings/asm/uci\_protocol.inc}}

\section{Rename And Copy A File}

\begin{listingbox}
\hspace*{2em}setptr ult\_buf, old\_name\
\hspace*{2em}setptr ult\_arg2, new\_name\
\hspace*{2em}jsr ultimate\_rename\
\
\hspace*{2em}setptr ult\_buf, new\_name\
\hspace*{2em}setptr ult\_arg2, copy\_name\
\hspace*{2em}jsr ultimate\_copy\
\
old\_name:  .byte "old.txt", 0\
new\_name:  .byte "new.txt", 0\
copy\_name: .byte "backup.txt", 0
\end{listingbox}

\begin{outputbox}
OLD.TXT becomes NEW.TXT; BACKUP.TXT is a copy of it.
\end{outputbox}

Both take the first name in \cmd{ult\_buf} and the second in \cmd{ult\_arg2}.
The request block carries two spans and these commands put three runs of bytes
on the wire, so the first name's own terminator separates it from the second.
The Ultimate performs the copy, so no bytes pass through the C64.

\evidence{\cmd{bindings/asm/ultimate.inc}; \cmd{src/uci/dosinfo.s}}

\section{Load And Save C64 Memory}

\begin{listingbox}
\hspace*{2em}setptr ult\_buf, prg\_name\
\hspace*{2em}setword ult\_addr, 0\
\hspace*{2em}jsr ultimate\_load\
\hspace*{2em}jsr ultimate\_last\_end\
\hspace*{2em}sta load\_end\
\hspace*{2em}stx load\_end + 1\
\
\hspace*{2em}setword ult\_addr, \$C000\
\hspace*{2em}jsr ultimate\_load\
\
\hspace*{2em}setptr ult\_buf, raw\_name\
\hspace*{2em}setword ult\_addr, \$C000\
\hspace*{2em}setword ult\_max, 2048\
\hspace*{2em}jsr ultimate\_bload\
\
\hspace*{2em}setptr ult\_buf, save\_name\
\hspace*{2em}setword ult\_addr, \$0400\
\hspace*{2em}setword ult\_max, 1000\
\hspace*{2em}jsr ultimate\_save\
\
prg\_name:  .byte "game.prg", 0\
raw\_name:  .byte "font.bin", 0\
save\_name: .byte "screen.bin", 0\
load\_end:  .res 2
\end{listingbox}

\begin{outputbox}
load_end = address after loaded PRG data\\
FONT.BIN loads at \$C000; SCREEN.BIN receives 1000 bytes.
\end{outputbox}

An \cmd{ult\_addr} of zero makes load use the PRG header; \cmd{\$C000} shows
the explicit-address form. \cmd{ultimate\_last\_end} returns the previous load
end in \cmd{A}/\cmd{X}. Bload consumes no header and obeys \cmd{ult\_max}; save
uses \cmd{ult\_max} as its length.

\evidence{\cmd{docs/asm-abi.md}; \cmd{src/uci/file.s}}

\section{Mount A Disk Image}

The drive is named by the number it answers to on the IEC serial bus, 8 or 9,
rather than by a slot. \cmd{ULTIMATE\_DRIVE\_LAST} is zero and asks for
whichever drive was mounted into last, or drive~A when nothing has been.

\begin{listingbox}
\hspace*{2em}setptr ult\_buf, image\_name\
\hspace*{2em}lda \#8\
\hspace*{2em}jsr ultimate\_mount\
\hspace*{2em}cmp \#ULTIMATE\_OK\
\hspace*{2em}bne failed\
\
\hspace*{2em}lda \#8\
\hspace*{2em}jsr ultimate\_swap\
\
\hspace*{2em}lda \#8\
\hspace*{2em}jsr ultimate\_unmount\
\
image\_name: .byte "games/pitfall.d64", 0
\end{listingbox}

\begin{outputbox}
PITFALL.D64 is mounted, swapped and unmounted; A = ULTIMATE_OK each time.
\end{outputbox}

The device number goes in \cmd{A} and is reloaded before each call, because
every entry point returns its result there. \cmd{ultimate\_mount} takes the
image name through \cmd{ult\_buf} and measures it itself.

The firmware picks the image type from the extension: \cmd{.D64}, \cmd{.D71}
and \cmd{.D81} hold the sectors of a 1541, 1571 and 1581 disk, and \cmd{.G64}
and \cmd{.G71} hold raw track data. Mounting into a drive that is switched off
returns \cmd{ULTIMATE\_ERR\_DEVICE}.

\evidence{\cmd{bindings/asm/ultimate.inc}; \cmd{src/uci/disk.s};
\cmd{docs/uci.md}}

\Needspace{30\baselineskip}
\section{Control The Machine And Its Drives}

\begin{listingbox}
\hspace*{2em}lda \#<drives\
\hspace*{2em}ldx \#>drives\
\hspace*{2em}jsr ultimate\_drive\_info\
\hspace*{2em}cmp \#ULTIMATE\_OK\
\hspace*{2em}bne failed\
\hspace*{2em}lda drives + CTRL\_DRVINFO\_COUNT\
\hspace*{2em}beq no\_drives\
\hspace*{2em}lda drives + CTRL\_DRVINFO\_FIRST + 1\
\hspace*{2em}sta first\_device\
\
\hspace*{2em}lda \#ULTIMATE\_DRIVE\_A\
\hspace*{2em}ldx \#1\
\hspace*{2em}jsr ultimate\_drive\_enable\
\
\hspace*{2em}lda \#ULTIMATE\_DRIVE\_A\
\hspace*{2em}jsr ultimate\_drive\_power\
\hspace*{2em}lda ult\_stage\
\hspace*{2em}beq drive\_is\_off\
\
drives:       .res CTRL\_DRVINFO\_BYTES\
first\_device: .res 1
\end{listingbox}

\begin{outputbox}
drives holds a count and one record per drive; ult_stage = 1 when it runs.
\end{outputbox}

\cmd{ultimate\_drive\_info} takes the block pointer in \cmd{A}/\cmd{X}. The
block is a count at \cmd{CTRL\_DRVINFO\_COUNT}, then
\cmd{CTRL\_DRVINFO\_RECORD} bytes per drive from \cmd{CTRL\_DRVINFO\_FIRST}:
type, then the IEC device number, then 1 while the drive is running. On any
failure the count is set to zero, so a caller that skips the result check reads
no records.

\cmd{ultimate\_drive\_enable} takes the drive in \cmd{A} and the state in
\cmd{X}: zero switches it off, and any non-zero value switches it on.
\cmd{ultimate\_drive\_power} takes the drive in \cmd{A} and reports the state
in \cmd{ult\_stage}, which is 1 while the drive runs.

\begin{note}
Two numbering schemes meet here. \cmd{ultimate\_drive\_enable} and
\cmd{ultimate\_drive\_power} take \cmd{ULTIMATE\_DRIVE\_A} or
\cmd{ULTIMATE\_DRIVE\_B}, which name a physical drive.
\cmd{ultimate\_mount} takes the device number the drive answers to on the IEC
bus. \cmd{ultimate\_drive\_info} connects the two: each record carries both.
\end{note}

\begin{warning}
\cmd{ultimate\_drive\_power} writes the flag to \cmd{ult\_stage} and uses
\cmd{ult\_stage + 1} to \cmd{ult\_stage + 3} as scratch. Read the flag before
the next SDK call.
\end{warning}

\subsection*{Reset the machine, or open its menu}

\begin{listingbox}
\hspace*{2em}jsr ultimate\_freeze\
\
\hspace*{2em}jsr ultimate\_reboot\
\hspace*{2em}; nothing here runs
\end{listingbox}

\begin{outputbox}
freeze returns when the menu is left; reboot does not return.
\end{outputbox}

\cmd{ultimate\_reboot} resets the C64, so the program calling it stops existing
part way through the call. \cmd{ultimate\_freeze} stops the C64 and shows the
Ultimate's menu, and returns when whoever is at the keyboard leaves it.

\subsection*{Read the RAM disk layout}

\begin{listingbox}
\hspace*{2em}lda \#<ramdisk\
\hspace*{2em}ldx \#>ramdisk\
\hspace*{2em}jsr ultimate\_ramdisk\_info\
\
ramdisk: .res CTRL\_RAMDISK\_BYTES
\end{listingbox}

\begin{outputbox}
ramdisk holds CTRL_RAMDISK_DRIVES two-byte records.
\end{outputbox}

Each record is a drive type code and a size in 64K units. A type of zero is a
slot with nothing in it. These are the RAM disks GEOS MegaPatch uses.

\evidence{\cmd{bindings/asm/ultimate.inc}; \cmd{src/uci/control.s};
\cmd{docs/uci.md}}

\Needspace{30\baselineskip}
\section{Read And Set The Clock}

\begin{listingbox}
\hspace*{2em}setptr ult\_buf, time\_text\
\hspace*{2em}setword ult\_buflen, ULTIMATE\_TIME\_BUFFER\
\hspace*{2em}setword ult\_outlen, 0\
\hspace*{2em}lda \#ULTIMATE\_TIME\_PLAIN\
\hspace*{2em}jsr ultimate\_get\_time\
\
time\_text: .res ULTIMATE\_TIME\_BUFFER
\end{listingbox}

\begin{outputbox}
time_text = "2026/08/22 14:30:00"
\end{outputbox}

\cmd{A} is the format. \cmd{ULTIMATE\_TIME\_PLAIN} asks for the date and time
and \cmd{ULTIMATE\_TIME\_WEEKDAY} puts the weekday in front of it.
\cmd{ULTIMATE\_TIME\_BUFFER} is a buffer size that fits either form. This is
one of the few entry points that reads \cmd{ult\_buflen} as an input rather
than computing it. A buffer too small still returns \cmd{ULTIMATE\_OK}, with
the part that fitted.

\subsection*{Set it}

\begin{listingbox}
\hspace*{2em}lda \#126\hspace{2em}; 2026, as the year less 1900\
\hspace*{2em}sta ult\_stage + 0\
\hspace*{2em}lda \#8\hspace{3em}; month\
\hspace*{2em}sta ult\_stage + 1\
\hspace*{2em}lda \#22\hspace{2em}; day\
\hspace*{2em}sta ult\_stage + 2\
\hspace*{2em}lda \#14\hspace{2em}; hour\
\hspace*{2em}sta ult\_stage + 3\
\hspace*{2em}lda \#30\hspace{2em}; minute\
\hspace*{2em}sta ult\_stage + 4\
\hspace*{2em}lda \#0\hspace{3em}; second\
\hspace*{2em}sta ult\_stage + 5\
\hspace*{2em}jsr ultimate\_set\_time
\end{listingbox}

\begin{outputbox}
The clock is set to 2026/08/22 14:30:00; A = ULTIMATE_OK.
\end{outputbox}

The six bytes go in \cmd{ult\_stage} in that order. \cmd{ult\_stagelen} is not
a caller input here; the command always sends six bytes.

\begin{warning}
\cmd{ult\_stage} is shared. \cmd{ultimate\_mount}, \cmd{ultimate\_unmount},
\cmd{ultimate\_swap}, \cmd{ultimate\_get\_time},
\cmd{ultimate\_drive\_info} and \cmd{ultimate\_http\_body} all overwrite
\cmd{ult\_stage + 0}, and \cmd{ultimate\_drive\_power} and the keyed body calls
overwrite more of it. Fill it immediately before \cmd{ultimate\_set\_time}.
\end{warning}

\evidence{\cmd{bindings/asm/ultimate.inc}; \cmd{src/uci/clock.s};
\cmd{docs/uci.md}}

\section{Move Bytes Through The REU}

\begin{listingbox}
\hspace*{2em}jsr ultimate\_reu\_available\
\hspace*{2em}beq no\_reu\
\hspace*{2em}jsr ultimate\_reu\_size\
\hspace*{2em}sta reu\_banks\
\hspace*{2em}stx reu\_banks + 1\
\
\hspace*{2em}setword ult\_addr, \$0400\
\hspace*{2em}lda \#0\
\hspace*{2em}sta ult\_reu\
\hspace*{2em}sta ult\_reu + 1\
\hspace*{2em}sta ult\_reu + 2\
\hspace*{2em}sta ult\_reu + 3\
\hspace*{2em}setword ult\_reulen, 1000\
\hspace*{2em}lda \#0\
\hspace*{2em}sta ult\_reulen + 2\
\hspace*{2em}sta ult\_reulen + 3\
\hspace*{2em}jsr ultimate\_reu\_stash\
\hspace*{2em}jsr ultimate\_reu\_fetch\
\
reu\_banks: .res 2
\end{listingbox}

\begin{outputbox}
reu_banks = installed REU size in 64K banks\\
1000 bytes stashed and fetched; A = ULTIMATE_OK.
\end{outputbox}

Size returns 64K banks in \cmd{A}/\cmd{X}. Stash copies C64 RAM to the REU;
fetch copies it back. Both read the C64 address, 32-bit REU address and 32-bit
length from the shared variables.

\subsection*{Move the current file directly to and from the REU}

Both transfers below use REU address zero and a 65536-byte length:

\begin{listingbox}
\hspace*{2em}lda \#0\
\hspace*{2em}sta ult\_reu\
\hspace*{2em}sta ult\_reu + 1\
\hspace*{2em}sta ult\_reu + 2\
\hspace*{2em}sta ult\_reu + 3\
\hspace*{2em}sta ult\_reulen\
\hspace*{2em}sta ult\_reulen + 1\
\hspace*{2em}lda \#1\
\hspace*{2em}sta ult\_reulen + 2\
\hspace*{2em}lda \#0\
\hspace*{2em}sta ult\_reulen + 3
\end{listingbox}

\begin{outputbox}
ult_reu = \$00000000; ult_reulen = 65536.
\end{outputbox}

\Needspace{25\baselineskip}
\subsection*{File to REU}

\begin{listingbox}
\hspace*{2em}; Open ASSETS.BIN first; it becomes the current DOS file.\
\hspace*{2em}setptr ult\_buf, asset\_name\
\hspace*{2em}lda \#DOS\_FA\_READ\
\hspace*{2em}jsr ultimate\_open\
\hspace*{2em}cmp \#ULTIMATE\_OK\
\hspace*{2em}bne failed\
\hspace*{2em}lda \#0\
\hspace*{2em}sta ult\_num\
\hspace*{2em}sta ult\_num + 1\
\hspace*{2em}sta ult\_num + 2\
\hspace*{2em}sta ult\_num + 3\
\hspace*{2em}jsr ultimate\_seek\
\hspace*{2em}jsr ultimate\_reu\_load\
\hspace*{2em}jsr ultimate\_close\
\
asset\_name: .byte "assets.bin", 0
\end{listingbox}

\begin{outputbox}
ASSETS.BIN to REU \$000000: 65536 bytes
\end{outputbox}

\Needspace{18\baselineskip}
\subsection*{REU to file}

\begin{listingbox}
\hspace*{2em}; Open a writable current file before the transfer.\
\hspace*{2em}setptr ult\_buf, capture\_name\
\hspace*{2em}lda \#DOS\_FA\_CREATE\_ALWAYS | DOS\_FA\_WRITE\
\hspace*{2em}jsr ultimate\_open\
\hspace*{2em}cmp \#ULTIMATE\_OK\
\hspace*{2em}bne failed\
\hspace*{2em}jsr ultimate\_reu\_save\
\hspace*{2em}jsr ultimate\_close\
\
capture\_name: .byte "capture.bin", 0
\end{listingbox}

\begin{outputbox}
REU \$000000 to CAPTURE.BIN: 65536 bytes
\end{outputbox}

The REU address is zero and the length bytes encode 65536. Load and save take
no filename: each operates on the firmware's current open file established by
the preceding \cmd{ultimate\_open}.

\evidence{\cmd{docs/asm-abi.md}; \cmd{src/uci/reu.s}}

\Needspace{32\baselineskip}
\section{Inspect Network Interfaces}

\begin{listingbox}
\hspace*{2em}jsr ultimate\_net\_ifcount\
\hspace*{2em}cmp \#ULTIMATE\_OK\
\hspace*{2em}bne failed\
\hspace*{2em}lda ult\_iface\
\hspace*{2em}beq no\_interfaces\
\
\hspace*{2em}lda \#0\
\hspace*{2em}sta ult\_iface\
\hspace*{2em}setptr ult\_buf, mac\
\hspace*{2em}jsr ultimate\_net\_macaddr\
\hspace*{2em}setptr ult\_buf, ipconfig\
\hspace*{2em}jsr ultimate\_net\_ipconfig\
\
mac:      .res 6\
ipconfig: .res 12
\end{listingbox}

\begin{outputbox}
ult_iface = interface count; mac receives 6 bytes; ipconfig receives 12.
\end{outputbox}

The count is returned in \cmd{ult\_iface}; the address calls take an interface
index there and place fixed-size replies at \cmd{ult\_buf}.

\subsection*{Change an interface's address}

\begin{listingbox}
\hspace*{2em}lda \#0\
\hspace*{2em}sta ult\_iface\
\hspace*{2em}setptr ult\_buf, ipconfig\
\hspace*{2em}jsr ultimate\_net\_ipconfig\
\hspace*{2em}cmp \#ULTIMATE\_OK\
\hspace*{2em}bne failed\
\
\hspace*{2em}lda \#60\
\hspace*{2em}sta ipconfig + 3\
\hspace*{2em}jsr ultimate\_net\_setip
\end{listingbox}

\begin{outputbox}
The interface keeps its netmask and gateway and answers on .60 instead.
\end{outputbox}

\cmd{ultimate\_net\_setip} sends back the same twelve bytes
\cmd{ultimate\_net\_ipconfig} returned: four of address, then four of netmask,
then four of gateway. It reads \cmd{ult\_iface} and \cmd{ult\_buf}, and always
sends twelve bytes, so \cmd{ult\_buflen} is not an input.

This changes the running configuration only; nothing is written to the
Ultimate's stored settings.

\section{Use TCP And UDP Sockets}

\begin{listingbox}
\hspace*{2em}setptr ult\_buf, host\
\hspace*{2em}setword ult\_port, 6502\
\hspace*{2em}jsr ultimate\_net\_connect\
\hspace*{2em}cmp \#ULTIMATE\_OK\
\hspace*{2em}bne failed\
\
\hspace*{2em}setptr ult\_buf, message\
\hspace*{2em}setword ult\_buflen, message\_end - message\
\hspace*{2em}jsr ultimate\_net\_write\
\hspace*{2em}; ult\_socklen is now the number accepted.\
\
host:        .byte "192.168.1.50", 0\
message:     .byte "hi", 13\
message\_end:
\end{listingbox}

\begin{outputbox}
TCP handle is in ult_sock; ult_socklen = 3 bytes accepted.
\end{outputbox}

\Needspace{25\baselineskip}
\subsection*{Read and close the TCP socket}

\begin{listingbox}
poll:\
\hspace*{2em}setptr ult\_buf, socket\_buffer\
\hspace*{2em}setword ult\_socklen, 128\
\hspace*{2em}jsr ultimate\_net\_read\
\hspace*{2em}cmp \#ULTIMATE\_END\
\hspace*{2em}beq peer\_closed\
\hspace*{2em}cmp \#ULTIMATE\_OK\
\hspace*{2em}bne close\_socket\
\hspace*{2em}lda ult\_socklen\
\hspace*{2em}ora ult\_socklen + 1\
\hspace*{2em}beq poll\
\
close\_socket:\
\hspace*{2em}jsr ultimate\_net\_close\
\
socket\_buffer: .res 128
\end{listingbox}

\begin{outputbox}
Reply bytes are in socket_buffer; close returns ULTIMATE_OK.
\end{outputbox}

\Needspace{14\baselineskip}
\subsection*{UDP}

\begin{listingbox}
\hspace*{2em}setptr ult\_buf, host\
\hspace*{2em}setword ult\_port, 6502\
\hspace*{2em}jsr ultimate\_net\_udp\
\hspace*{2em}cmp \#ULTIMATE\_OK\
\hspace*{2em}bne failed\
\hspace*{2em}jsr ultimate\_net\_close
\end{listingbox}

\begin{outputbox}
UDP handle opens and closes; A = ULTIMATE_OK.
\end{outputbox}

Connect and UDP return the handle in \cmd{ult\_sock}. Write takes its length
from \cmd{ult\_buflen} and reports accepted bytes in \cmd{ult\_socklen}. Read
takes capacity and returns stored bytes through \cmd{ult\_socklen}. A zero-byte
success means poll again. \cmd{ULTIMATE\_END} means the handle is already gone;
do not close that ended handle.

\evidence{\cmd{docs/asm-abi.md}; \cmd{src/uci/net.s}}

\section{Use The HTTP Client}

\begin{listingbox}
\hspace*{2em}setptr ult\_url, url\
\hspace*{2em}setptr ult\_buf, http\_reply\
\hspace*{2em}setword ult\_buflen, 512\
\hspace*{2em}jsr ultimate\_http\_get\
\hspace*{2em}; ult\_httplen is the number stored.\
\
url:           .byte "http://192.168.1.50/", 0\
accept\_header: .byte "accept: text/plain", 0\
http\_reply:    .res 512
\end{listingbox}

\begin{outputbox}
One-call GET stores response bytes; A = ULTIMATE_OK.
\end{outputbox}

\Needspace{26\baselineskip}
\subsection*{Multi-step GET without a body}

\begin{listingbox}
\hspace*{2em}setptr ult\_url, url\
\hspace*{2em}lda \#HTTP\_VERB\_GET\
\hspace*{2em}jsr ultimate\_http\_open\
\hspace*{2em}cmp \#ULTIMATE\_OK\
\hspace*{2em}bne failed\
\
\hspace*{2em}setptr ult\_url, accept\_header\
\hspace*{2em}jsr ultimate\_http\_header\
\hspace*{2em}lda \#HTTP\_BODY\_NONE\
\hspace*{2em}sta ult\_body\
\hspace*{2em}setptr ult\_buf, http\_reply\
\hspace*{2em}setword ult\_buflen, 512\
\hspace*{2em}jsr ultimate\_http\_exchange\
\hspace*{2em}jsr ultimate\_http\_close\
\
\end{listingbox}

\begin{outputbox}
GET without a body stores response bytes; close returns ULTIMATE_OK.
\end{outputbox}

\Needspace{22\baselineskip}
The one-call GET creates and releases its own slot. The multi-step form stores
the request handle in \cmd{ult\_http}, and header, exchange and close all read
it there. \cmd{ult\_body} holds the body handle, and
\cmd{HTTP\_BODY\_NONE} in it means no body.
\cmd{ultimate\_http\_free\_all} recovers slots abandoned by an earlier
program.

The next section creates a body and posts it.

\evidence{\cmd{docs/asm-abi.md}; \cmd{src/uci/http.s}}

\Needspace{32\baselineskip}
\section{Build A Request Body}

The body is built inside the Ultimate: create a slot, add to it one call at a
time, then run the exchange. The handle lives in \cmd{ult\_body}, where
\cmd{ultimate\_http\_exchange} already looks for it, so nothing has to carry it
around.

\begin{listingbox}
\hspace*{2em}lda \#HTTP\_BODY\_JSON\_OBJECT\
\hspace*{2em}jsr ultimate\_http\_body\
\hspace*{2em}cmp \#ULTIMATE\_OK\
\hspace*{2em}bne failed\
\
\hspace*{2em}setptr ult\_url, key\_name\
\hspace*{2em}setptr ult\_buf, value\_ada\
\hspace*{2em}jsr ultimate\_http\_body\_string\
\
\hspace*{2em}setptr ult\_url, key\_score\
\hspace*{2em}setword ult\_val, 4200\
\hspace*{2em}lda \#0\
\hspace*{2em}sta ult\_val + 2\
\hspace*{2em}sta ult\_val + 3\
\hspace*{2em}jsr ultimate\_http\_body\_int\
\
key\_name:  .byte "name", 0\
key\_score: .byte "score", 0\
value\_ada: .byte "ada", 0
\end{listingbox}

\begin{outputbox}
The body holds \{"NAME":"ADA","SCORE":4200\}.
\end{outputbox}

\cmd{ult\_url} is the key for every keyed call, and it is the same variable
\cmd{ultimate\_http\_open} uses for a URL. \cmd{ult\_buf} is a string value.
\cmd{ult\_val} is four bytes, low byte first, and all four are sent, because
the firmware sign-extends from the last byte it receives.

\cmd{ultimate\_http\_body\_bool} reads \cmd{ult\_val + 0} only, and any
non-zero value is true.

\subsection*{Send it}

\begin{listingbox}
\hspace*{2em}setptr ult\_url, url\
\hspace*{2em}lda \#HTTP\_VERB\_POST\
\hspace*{2em}jsr ultimate\_http\_open\
\hspace*{2em}cmp \#ULTIMATE\_OK\
\hspace*{2em}bne failed\
\
\hspace*{2em}setptr ult\_buf, http\_reply\
\hspace*{2em}setword ult\_buflen, 512\
\hspace*{2em}jsr ultimate\_http\_exchange\
\hspace*{2em}jsr ultimate\_http\_close\
\hspace*{2em}jsr ultimate\_http\_body\_free
\end{listingbox}

\begin{outputbox}
The response is in http_reply; ult_httplen is the number stored.
\end{outputbox}

Set \cmd{ult\_buf} back to the reply buffer before the exchange. The keyed body
calls used it for the last string value, and the exchange uses it for the
reply.

\subsection*{Nest an object or an array}

\begin{listingbox}
\hspace*{2em}setptr ult\_url, key\_stats\
\hspace*{2em}jsr ultimate\_http\_body\_object\
\hspace*{2em}; keys added here go inside "stats"\
\hspace*{2em}jsr ultimate\_http\_body\_up\
\
\hspace*{2em}setptr ult\_url, key\_runs\
\hspace*{2em}jsr ultimate\_http\_body\_array\
\hspace*{2em}; inside an array the key is ignored\
\hspace*{2em}jsr ultimate\_http\_body\_up
\end{listingbox}

\begin{outputbox}
\{"STATS":\{...\},"RUNS":[...]\}
\end{outputbox}

Adding an object or an array enters it. Every key added afterwards goes inside,
until \cmd{ultimate\_http\_body\_up} steps back out to the parent.

\subsection*{Send raw bytes instead}

\begin{listingbox}
\hspace*{2em}lda \#HTTP\_BODY\_BINARY\
\hspace*{2em}jsr ultimate\_http\_body\
\
\hspace*{2em}setptr ult\_buf, block\
\hspace*{2em}setword ult\_buflen, block\_end - block\
\hspace*{2em}jsr ultimate\_http\_body\_binary\
\
\hspace*{2em}jsr ultimate\_http\_body\_clear\
\hspace*{2em}jsr ultimate\_http\_body\_free\
\
block:     .byte 1, 2, 3, 4\
block\_end:
\end{listingbox}

\begin{outputbox}
The body holds four bytes, then none; the slot is returned.
\end{outputbox}

\cmd{ultimate\_http\_body\_binary} is the other entry that reads
\cmd{ult\_buflen} as an input. Successive calls append.
\cmd{ultimate\_http\_body\_clear} empties a body and keeps its slot and format;
\cmd{ultimate\_http\_body\_free} returns the slot.

A \cmd{HTTP\_BODY\_BINARY} body takes \cmd{ultimate\_http\_body\_binary} and
nothing else; the JSON and form formats take everything else. A key is staged
through \cmd{ult\_stage}, so it is limited to \cmd{ULTIMATE\_HTTP\_KEY\_MAX}
bytes, which is 34. A longer key returns
\cmd{ULTIMATE\_ERR\_INVALID\_ARGUMENT} and never reaches the wire.

\begin{warning}
The firmware has sixteen slots for the whole machine, and a program that
crashes returns none of them. Free every body and every request.
\cmd{ultimate\_http\_free\_all} takes back headers and bodies together.
\end{warning}

\evidence{\cmd{bindings/asm/ultimate.inc}; \cmd{src/uci/httpbody.s}}

\section{Probe And Configure The Raw Transport}

\begin{listingbox}
\hspace*{2em}jsr uci\_present\
\hspace*{2em}beq no\_interface\
\hspace*{2em}jsr uci\_ident\
\hspace*{2em}sta raw\_ident\
\
\hspace*{2em}jsr uci\_get\_timeout\_a\
\hspace*{2em}pha\
\hspace*{2em}lda \#UCI\_TIMEOUT\_FOREVER\
\hspace*{2em}jsr uci\_set\_timeout\_a\
\hspace*{2em}jsr run\_slow\_command\
\hspace*{2em}tax\
\hspace*{2em}pla\
\hspace*{2em}jsr uci\_set\_timeout\_a\
\hspace*{2em}txa\
\
\hspace*{2em}jsr uci\_last\_code\
\hspace*{2em}sta device\_code\
\hspace*{2em}stx device\_code + 1\
\hspace*{2em}jsr uci\_abort
\end{listingbox}

\begin{outputbox}
uci_present returns nonzero; raw_ident receives the interface ID.\\
device_code receives the last target status; abort returns ULTIMATE_OK.
\end{outputbox}

Present and ident are direct register reads. Timeout zero waits forever; this
pattern restores the previous budget while preserving the command result.
Last code returns the raw target status in \cmd{A}/\cmd{X}. Abort releases an
unfinished exchange and is safe while idle.

\section{Raw Commands: Optional Arguments}

\begin{listingbox}
\hspace*{2em}jsr uci\_req\_clear\
\hspace*{2em}lda \#UCI\_TARGET\_DOS1\
\hspace*{2em}sta uci\_req + UCI\_REQ\_TARGET\
\hspace*{2em}lda \#UCI\_CMD\_IDENTIFY\
\hspace*{2em}sta uci\_req + UCI\_REQ\_COMMAND\
\hspace*{2em}setptr uci\_req + UCI\_REQ\_DATA, raw\_reply\
\hspace*{2em}setword uci\_req + UCI\_REQ\_DATAMAX, 64\
\hspace*{2em}jsr uci\_exec\_block\
\hspace*{2em}jsr uci\_req\_clear\
\hspace*{2em}lda \#UCI\_TARGET\_DOS1\
\hspace*{2em}sta uci\_req + UCI\_REQ\_TARGET\
\hspace*{2em}lda \#DOS\_CMD\_GET\_TIME\
\hspace*{2em}sta uci\_req + UCI\_REQ\_COMMAND\
\hspace*{2em}setptr uci\_req + UCI\_REQ\_ARGS, time\_format\
\hspace*{2em}setword uci\_req + UCI\_REQ\_ARGLEN, 1\
\hspace*{2em}jsr uci\_exec\_block\
\hspace*{2em}jsr uci\_req\_clear\
\hspace*{2em}lda \#UCI\_TARGET\_CONTROL\
\hspace*{2em}sta uci\_req + UCI\_REQ\_TARGET\
\hspace*{2em}lda \#CTRL\_CMD\_GET\_HWINFO\
\hspace*{2em}sta uci\_req + UCI\_REQ\_COMMAND\
\hspace*{2em}setptr uci\_req + UCI\_REQ\_DATA, raw\_reply\
\hspace*{2em}setword uci\_req + UCI\_REQ\_DATAMAX, 64\
\hspace*{2em}jsr uci\_exec\_block\
\hspace*{2em}setptr uci\_req + UCI\_REQ\_ARGS, hwinfo\_model\
\hspace*{2em}setword uci\_req + UCI\_REQ\_ARGLEN, 1\
\hspace*{2em}jsr uci\_exec\_block\
time\_format: .byte 0\
hwinfo\_model: .byte CTRL\_HWINFO\_MODEL\
raw\_reply:   .res 64
\end{listingbox}

\begin{outputbox}
IDENTIFY returns the target name.\\
GET_TIME returns the formatted time.\\
Both GET_HWINFO forms return the model name.
\end{outputbox}

Clearing the request leaves every optional span empty. The time command then
supplies its optional format byte. \cmd{GET\_HWINFO} is sent without and then
with its model selector, which is what an omitted optional argument means on
the wire. Set other spans only when needed.

\subsection*{Use a caller-owned request block}

\begin{listingbox}
\hspace*{2em}lda \#<my\_request\
\hspace*{2em}ldx \#>my\_request\
\hspace*{2em}jsr uci\_exec\
\
my\_request: .res UCI\_REQ\_SIZE
\end{listingbox}

\begin{outputbox}
A = request result; lengths and retained bytes are written into my_request.
\end{outputbox}

\cmd{uci\_exec} takes the request pointer directly. Use it instead of the
global block when the program owns more than one request.

\Needspace{30\baselineskip}
\section{Read A Raw Multi-block Reply}

\begin{listingbox}
\hspace*{2em}lda \#<my\_request\
\hspace*{2em}ldx \#>my\_request\
\hspace*{2em}jsr uci\_exec\_first\
\hspace*{2em}cmp \#ULTIMATE\_OK\
\hspace*{2em}bne failed\
\
next\_block:\
\hspace*{2em}; use this block's UCI\_REQ\_DATALEN bytes\
\hspace*{2em}lda uci\_more\
\hspace*{2em}beq all\_blocks\
\hspace*{2em}jsr uci\_exec\_next\
\hspace*{2em}cmp \#ULTIMATE\_OK\
\hspace*{2em}beq next\_block
\end{listingbox}

\begin{outputbox}
Each iteration receives one block; uci_more = 0 after the final block.
\end{outputbox}

First submits the request and returns one block. \cmd{uci\_more} is the boolean
state byte that says whether another follows; next releases the current block
and receives it. Issue no unrelated command until the final block or call
\cmd{uci\_abort} when stopping early.

\evidence{\cmd{docs/asm-abi.md}; \cmd{src/uci/uci\_core.s}}
